\subsection*{Applicazione della teoria dei momenti generalizzati allo spin}
Non c'è molto da dire: siccome lo spin è il momento angolare intrinseco e rispetta la proprietà essenziale dei momenti angolari, si procede senza variare nulla se non che ora $\hat J$ è $\S$; tuttavia è essenziale notare che le affermazioni fatte nel capitolo precedente non sono valide, in quanto lo spazio ambiene era $\R^3$, e per di più si era supposta l'invarianza per rotazioni di $2\pi$ (tendenzialmente non vera nel caso dello spin).

% Sì l'intera lezione era su questo.

Risulta che l'evoluzione temporale di una sovrapposizione di stati di momento angolare è periodica: si parla della \textbf{Precessione di spin}\index{Precessione di spin}.